\section{Appendix: Hilbert-Poincar\'e series and Gr\"obner bases}\label{appendix}
In this section we collect some of the basics about the theory of
Hilbert-Poincar\'{e} series. For a detailed introduction, especially
proofs, we refer to \cite{greuel_pfister}. We also recall some results on Gr\"obner basis theory from \cite{Cox_Little_Oshea}.\\

Let $A$ be a ($\Z$-)graded $k$-algebra and let $M=\oplus_{i\in \Z} M_i$ be a
graded $A$-module with $i$th graded pieces $A_i$ and $M_i$ of finite
$k$-dimension. The {\it Hilbert function $H_M\colon \Z \ra \Z$ of $M$}
is defined by $H_M(i)= \dim_k M_i$, and its corresponding generating
series $$\HP_M(t)=\sum_{i\in \Z}H_M(i)t^i \in \Z((t))$$ is
called the {\it Hilbert-Poincar\'{e} series of $M$}. It is well-known
that if $A$ is a Noetherian $k$-algebra generated by homogeneous
elements $x_1, \ldots, x_n$ of degrees $d_1,\ldots, d_n$ and $M$ is a finitely
generated $A$-module
then $$\HP_M(t)=\frac{Q_M(t)}{\prod_{i=1}^n(1-t^{d_i})}$$ for some
$Q_M(t) \in \Z[t]$ which is called the {\it (weighted) first Hilbert
 series of $M$}. If $A$ respectively $M$ is non-Noetherian then the
Hilbert-Poincar\'{e} series of $M$ need not be rational anymore. For
the rest of this section we assume that the polynomial ring
$k[x_1,\ldots, x_n]$ is graded (not necessarily standard graded). The
notions of homogeneous ideal and degree are to be understood
relative to this grading. If $M$ is graded then for any integer $d$ we
write $M(d)$ for the {\it $d$th twist of $M$}, i.e., the graded
$A$-module with $M(d)_i=M_{i+d}$.\\

The following Lemma follows immediately from additivity of dimension:  

\begin{samepage}
\begin{Lem}[Lemma 5.1.2 in \cite{greuel_pfister}]\label{lem:hilbert_trick_1}
  Let $A$ and $M$ be as above. Let $d$ be a non-negative integer,
  $f\in A_d$ and $\varphi\colon M(-d)\ra M$ be defined by $\varphi(m)
  = f\cdot m$; then $\ker(\varphi)$ and $\coker(\varphi)$ are graded
  $A/(f)$-modules with the induced gradings and $$\HP_M(t) = t^d\cdot
  \HP_M(t) + \HP_{\coker(\varphi)}(t) - t^d\cdot \HP_{\ker(\varphi)}(t).$$
\end{Lem}
\end{samepage}

As an immediate consequence we obtain the useful:

\begin{samepage}
\begin{Cor}[Lemma 5.2.2 in
  \cite{greuel_pfister}]\label{cor:hilbert_hilbert_trick_2}
  Let $I\subseteq k[x_1,\ldots, x_n]$ be a homogeneous ideal, and let
  $f\in k[x_1,\ldots, x_n]$ be a homogeneous polynomial of degree $d$
  then $$\HP_{k[x]/I}(t) = \HP_{k[x]/(I,f)}(t) + t^d\HP_{k[x]/(I:f)}(t).$$
\end{Cor}
\end{samepage}

For homogeneous ideals the leading ideal already determines the
Hilbert-Poincar\'{e} series. After fixing a monomial order, the \textit{leading ideal} $L(I)$ of an ideal $I$ in $k[x_1,\ldots,x_n]$ is defined as the (monomial) ideal generated by the leading monomials of all elements in $I$. Then one has:

\begin{samepage}
\begin{Thm}[Theorem 5.2.6 in
  \cite{greuel_pfister}]\label{thm:hilbert_leading} Let $>$ be any
  monomial ordering on $k[x_1,\ldots, x_n]$, let $I\subseteq k[x]$
  be a homogeneous ideal and denote by $L(I)$ its leading ideal with
  respect to $>$. Then $$\HP_{k[x]/I}(t) = \HP_{k[x]/L(I)}(t).$$
\end{Thm}
\end{samepage}

To compute the leading ideal one can use Gr\"obner bases. Let $I$ be an ideal in the polynomial ring $k[x_1,\ldots, x_n]$ with a fixed monomial order $<$. Then $\{g_1,\ldots, g_l\}\subset I$ is called a \textit{Gr\"obner basis} of $I$ if $L(I)$ is generated by $\{ \lm(g_i) \,;\, 1\leq i\leq l\}$, where we write $\lm$ for `leading monomial'. For $f\in k[x_1,\ldots, x_n]$ and a subset $H = \{h_1,\ldots, h_s\}$ of $k[x_1,\ldots ,x_n]$ one says that $f$ \textit{reduces to zero modulo} $H$ if 
\[  f = a_1h_1 + \cdots + a_sh_s \]
for $a_i \in k[x_1,\ldots, x_n]$, such that $\lm(f) \geq \lm(a_ih_i) $ whenever $a_ih_i\neq 0$. One writes $f \rightarrow_H 0$.

Finally we need the definition of a syzygy. Let $\mathcal{F}=(f_1,\ldots ,f_s)\in (k[x_1,\ldots,x_n])^s$. A \textit{syzygy} on the leading terms of the $f_i$ is an $s$-tuple $(h_1,\ldots,h_s) \in (k[x_1,\ldots,x_n])^s$ such that 
\[  \sum_{i=1}^s h_i\, \lt(f_i) = 0,\]
where $\lt$ stands for `leading term'. The set of syzygies $S(\mathcal{F})$ on the leading terms of $\mathcal{F}$ form a $k[x_1,\ldots,x_n]$-submodule of $(k[x_1,\ldots,x_n])^s$. A generating set of this module is called a \textit{basis}. With $F=\{f_1,\ldots, f_s\}$, we will say that a syzygy $(h_1,\ldots, h_s)\in S(\mathcal{F})$ \textit{reduces to zero modulo} $F$ if
\[  \sum_{i=1}^s h_if_i \rightarrow_F 0.  \]
If each $h_i$ consists of a single term $c_ix^{\alpha_i}$ and $x^{\alpha_i}\lm(f_i)$ is a fixed monomial $x^{\alpha}$ if $c_i\neq 0$, then the syzygy $(h_1,\ldots, h_s)$ is called \textit{homogeneous of multidegree} $\alpha$. For $i<j$ let $x^{\gamma}$ be the least common multiple of the leading monomials of $f_i$ and $f_j$. One calls
\[ S(f_i,f_j) := \frac{x^{\gamma}}{\lt(f_i)}\, f_i - \frac{x^{\gamma}}{\lt(f_j)}\, f_j \] 
the \textit{$S$-polynomial} of $f_i$ and $f_j$. It gives rise to the homogeneous syzygy
\[  S_{i,j} := \frac{x^{\gamma}}{\lt(f_i)}\, e_i - \frac{x^{\gamma}}{\lt(f_j)}\, e_j, \]
where $e_i$ and $e_j$ denote standard basis vectors of $(k[x_1,\ldots,x_n])^s$. Then we have the following results:

\begin{Prop}[Proposition 4 p.103 in \cite{Cox_Little_Oshea}]\label{prop:coprime_leading_monomials} Let $G\subset k[x_1,\ldots,x_n]$ be a finite set. Assume that $f,g\in G$ have relatively prime leading monomials. Then $S(f,g)\rightarrow_G 0$.
\end{Prop}


\begin{Prop}[Proposition 8 p.105 in
  \cite{Cox_Little_Oshea}]\label{prop:S_polynomials_form_basis} For an $s$-tuple of polynomials $(f_1,\ldots ,f_s)\in (k[x_1,\ldots,x_n])^s$ we have that the set of all $S_{i,j}$ form a homogeneous basis of the syzygies on the leading terms of the $f_i$.  
\end{Prop}



\begin{Thm}[Theorem 9 p.106 in
  \cite{Cox_Little_Oshea}]\label{thm:criterion_for_Groebner_basis} Let $\mathcal{G}=(g_1,\ldots,g_s)$ be an $s$-tuple of polynomials and let $I$ be the ideal of $k[x_1,\ldots,x_n]$ generated by $G=\{g_1,\ldots,g_s\}$. Then $G$ is a Gr\"obner basis for $I$ if and only if every element of a homogeneous basis for the syzygies $S(\mathcal{G})$ reduces to zero modulo $G$.
\end{Thm}

\begin{Prop}[Proposition 10 p.107 in
  \cite{Cox_Little_Oshea}]\label{prop:reduction_of_syzygy_basis} Let $\mathcal{G}=(g_1,\ldots,g_s)$ be an $s$-tuple of polynomials. Suppose that we have a subset $S \subset \{S_{i,j} \,;\, 1\leq i <j \leq s\}$ that is a basis of $S(\mathcal{G})$. Moreover, suppose that we have distinct elements $g_i,g_j,g_k$ such that $\lm(g_k)$ divides the least common multiple of $\lm(g_i)$ and $\lm(g_j)$. If $S_{i,k},S_{j,k}\in S$, then $S\setminus \{S_{i,j}\}$ is also a basis of $S(\mathcal{G})$. Here we put $S_{i,j} := S_{j,i}$ if $i> j$. 
\end{Prop}
